\begin{proof}[Proof of \cref{obj:thm:oracle-frontier-resolved}]
Fix \(v\in\mathcal V\). Admissibility gives
\[
1<p\le2,\qquad \frac14\le\gamma\le1,\qquad
0<q=\frac{p-1}{p}\le\frac12,\qquad
2\gamma+q>0,\qquad E_0>0.
\]
We use the constants and scales in the statement.

\begin{enumerate}
\item Since \(-1\le-\gamma\) and the base \(4\) exceeds one,
\[
4^{-1}\le4^{-\gamma},
\qquad
\frac{1}{64\sqrt3}
=\frac{4^{-1}}{16\sqrt3}
\le c^{\mathrm e}_v.
\]
Also,
\[
C^{\mathrm e}_v
=3\left(120+128\sqrt{160\sqrt2}\right)>0.
\]
% lean: CausalSmith.Stat.FinitepHomogeneityDensegamma.cOracle_uniform_lower
% lean: CausalSmith.Stat.FinitepHomogeneityDensegamma.COr_pos

\item Fix an integer \(n\ge2\), and define the lower separation by
\[
\varrho_-=c^{\mathrm e}_v\rho_n^{\mathrm e}(v).
\]
By \cref{obj:lem:paired-tent-full-record-lower}, there are a null law \(P_0\in H_0(v)\) and a finite prior
\[
\pi_1=\sum_{j=1}^{k}\omega_j\delta_{P_j},
\qquad
\omega_j\ge0,\qquad
\sum_{j=1}^{k}\omega_j=1,
\qquad
\exists j:\omega_j>0.
\]
For every positive-weight support law, that result supplies
\[
P_j\in\mathcal M_v,\qquad
e_{P_j}(x)=\frac12=e_{P_0}(x)\quad(x\in[0,1]),
\qquad
\varrho_-\le d(P_j)<d_0\le D_v.
\]
It also supplies
\[
\mathbb P_{\pi_1,n}
=\sum_{j=1}^{k}\omega_jP_j^{\otimes n},
\qquad
\operatorname{TV}\!\left(P_0^{\otimes n},\mathbb P_{\pi_1,n}\right)<\frac12,
\qquad
B_n^{\mathrm e}(v,\varrho_-)\ge\frac25.
\]
% lean: CausalSmith.Stat.FinitepHomogeneityDensegamma.paired_tent_full_record_lower

Choose an index \(j_+\in\{1,\ldots,k\}\) with \(\omega_{j_+}>0\). Positivity of the lower multiplier and of \(n^{-E_0}\) yields
\[
0<\varrho_-\le d(P_{j_+})<D_v.
\]
Thus the alternatives at separation \(\varrho_-\) are nonempty. For every \(0<\varrho\le\varrho_-\), their inclusion in the alternatives at separation \(\varrho\) gives
\[
B_n^{\mathrm e}(v,\varrho_-)
\le B_n^{\mathrm e}(v,\varrho).
\]
The size constraint is the same in both infima. Consequently, every successful separation below \(D_v\) in \cref{obj:def:oracle-radius} is at least \(\varrho_-\), because
\[
B_n^{\mathrm e}(v,\varrho)\ge\frac25>\frac1{10}
\qquad(0<\varrho\le\varrho_-).
\]
The additional candidate \(D_v\) also exceeds \(\varrho_-\). Taking the infimum therefore gives
\[
c^{\mathrm e}_v\rho_n^{\mathrm e}(v)
=\varrho_-\le r_n^{*,\mathrm e}(v).
\]
% lean: CausalSmith.Stat.FinitepHomogeneityDensegamma.oracleCriticalRadius_ge_of_failed

\item We next establish the null calibration of the rule in \cref{obj:def:oracle-handle}. Use its tuning quantities
\[
s=\lfloor n/2\rfloor,\qquad a=s^{-E_0},\qquad
T=2^{\lceil\log_2(a^{-1/(p-1)})\rceil},\qquad
J=2^{\lceil\log_2(a^{-1/\gamma})\rceil},
\]
\[
b=20T^{1-p},\qquad
\Lambda=\frac{80T^{2-p}}s.
\]
Here \(s\ge1\), \(0<a\le1\), \(T\ge1\), \(J\ge1\), \(b\ge0\), and \(\Lambda>0\).
With \(F_J\) from \cref{obj:def:projection}, define
\[
u=J^{-1/2}(1,\ldots,1)^T,\qquad
C_J=I_J-uu^T.
\]
The cells have Lebesgue measure \(1/J\), so
\[
\int_0^1F_J(x)F_J(x)^T\,d\lambda(x)=I_J,
\qquad
\int_0^1F_J(x)\,d\lambda(x)=u,
\qquad
C_JF_J(x)=F_J(x)-u.
\]
Thus \(C_J\) is a symmetric orthogonal projection and
\[
\|C_J\xi_{\mathrm{dir}}\|_2\le\|\xi_{\mathrm{dir}}\|_2
\qquad(\xi_{\mathrm{dir}}\in\mathbb R^J).
\]

For \(P\in\mathcal M_v\), define the clipped conditional mean and second moment on \([0,1]\) by
\[
g_{T,P}(x)
=\int_{\mathbb R}\ell_T(y)\,Q_{1,P}(dy\mid x)
-\int_{\mathbb R}\ell_T(y)\,Q_{0,P}(dy\mid x),
\]
\[
m_{2,T,P}(x)
=\frac{\int_{\mathbb R}\ell_T(y)^2\,Q_{1,P}(dy\mid x)}{e_P(x)}
+\frac{\int_{\mathbb R}\ell_T(y)^2\,Q_{0,P}(dy\mid x)}{1-e_P(x)}.
\]
These functions are finite and measurable, since clipping is bounded and overlap makes both denominators positive.

For every \(y\in\mathbb R\),
\[
|y-\ell_T(y)|\le |y|^pT^{1-p},
\qquad
\ell_T(y)^2\le |y|^pT^{2-p}.
\]
Indeed, on \(|y|\le T\), the clipping remainder vanishes and
\(|y|^{2-p}\le T^{2-p}\); on \(|y|>T\), the remainder is at most \(|y|\) and the clipped square equals \(T^2\). The moment restriction in \cref{obj:ass:raw-moment} therefore gives, for each \(a'\in\{0,1\}\) and almost every \(x\),
\[
\left|\int_{\mathbb R}(y-\ell_T(y))\,Q_{a',P}(dy\mid x)\right|
\le10T^{1-p},
\qquad
\int_{\mathbb R}\ell_T(y)^2\,Q_{a',P}(dy\mid x)
\le10T^{2-p}.
\]
The same moment restriction ensures integrability of \(y\), using
\(|y|\le1+|y|^p\).
The arm-mean identities now yield
\[
g_{T,P}(x)-\tau_P(x)
=
\int_{\mathbb R}(y-\ell_T(y))\,Q_{0,P}(dy\mid x)
-\int_{\mathbb R}(y-\ell_T(y))\,Q_{1,P}(dy\mid x),
\]
and hence
\[
|g_{T,P}(x)-\tau_P(x)|\le b,
\qquad
m_{2,T,P}(x)
\le10T^{2-p}
\left(\frac1{e_P(x)}+\frac1{1-e_P(x)}\right)
\le80T^{2-p}
\]
almost everywhere. These are respectively the conditional mean and second moment bounds for \(R_i\), since the supplied true propensity agrees with its overlap-clipped version.
% lean: CausalSmith.Stat.FinitepHomogeneityDensegamma.oracle_clipping_moments

Define the common block mean and the two histogram coefficient vectors by
\[
\mu_P=\mathbb E_{P^{\otimes n}}[V_1]
=\mathbb E_{P^{\otimes n}}[V_2]
=C_J\theta_{T,P},
\]
\[
\theta_{T,P}=\int_0^1g_{T,P}(x)F_J(x)\,d\lambda(x),
\qquad
\theta_{\tau,P}=\int_0^1\tau_P(x)F_J(x)\,d\lambda(x).
\]
The equality of the block means follows by averaging \(s\) identically distributed records in each block.

For a null law, define its constant effect by
\[
c_P=\tau_P(0),\qquad \tau_P(x)=c_P\quad(x\in[0,1]).
\]
Because \(C_Ju=0\),
\[
\mu_P=C_J\int_0^1(g_{T,P}(x)-c_P)F_J(x)\,d\lambda(x).
\]
For every histogram coordinate \(j\),
\[
\left|
\left(\int_0^1(g_{T,P}-c_P)F_J\,d\lambda\right)_j
\right|
=
\left|\sqrt J\int_{I_{j,J}}(g_{T,P}-c_P)\,d\lambda\right|
\le\frac b{\sqrt J}.
\]
Summing the squared coordinate bounds and applying the contraction \(C_J\) gives
\[
\|\mu_P\|_2^2
\le\sum_{j=1}^{J}\frac{b^2}{J}=b^2,
\qquad
\|\mu_P\|_2\le b.
\]
% lean: CausalSmith.Stat.FinitepHomogeneityDensegamma.oracle_null_population_norm_le

For any model law and any \(\xi_{\mathrm{dir}}\in\mathbb R^J\), the conditional second moment bound and the feature isometry give
\[
\mathbb E_P\!\left[
\left\langle\xi_{\mathrm{dir}},C_JF_J(X)R\right\rangle^2
\right]
\le80T^{2-p}\int_0^1
\left\langle C_J\xi_{\mathrm{dir}},F_J(x)\right\rangle^2\,d\lambda(x)
\le80T^{2-p}\|\xi_{\mathrm{dir}}\|_2^2.
\]
Here the single-record score is defined by
\[
R=\frac{A\ell_T(Y)}{e_P(X)}
-\frac{(1-A)\ell_T(Y)}{1-e_P(X)}.
\]
Independence within a block gives, for \(b'\in\{1,2\}\),
\[
\operatorname{Var}_{P^{\otimes n}}
\!\left(\left\langle\xi_{\mathrm{dir}},V_{b'}\right\rangle\right)
=
\frac1s\operatorname{Var}_P
\!\left(\left\langle\xi_{\mathrm{dir}},C_JF_J(X)R\right\rangle\right)
\le\Lambda\|\xi_{\mathrm{dir}}\|_2^2.
\]
% lean: CausalSmith.Stat.FinitepHomogeneityDensegamma.oracleBlockVec_variance_le

Define
\[
\Delta_{b',P}=V_{b'}-\mu_P\quad(b'\in\{1,2\}),
\qquad
W_{\mathrm e}=\langle V_1,V_2\rangle.
\]
The centered vectors are independent and have mean zero. In particular,
\[
\mathbb E_{P^{\otimes n}}[W_{\mathrm e}]=\|\mu_P\|_2^2,
\]
\[
W_{\mathrm e}-\|\mu_P\|_2^2
=
\langle\mu_P,\Delta_{1,P}\rangle
+\langle\mu_P,\Delta_{2,P}\rangle
+\langle\Delta_{1,P},\Delta_{2,P}\rangle.
\]
The expectations of all three cross products in the square of this sum vanish: the product of the two linear terms factors into their zero means, and conditioning on the other block makes each linear--bilinear cross product zero. To bound the remaining bilinear square, let
\[
\mathbf e^{(j)}
=\bigl(\mathbf1_{\{j'=j\}}\bigr)_{j'=1}^{J}
\qquad(j=1,\ldots,J)
\]
be the coordinate unit vectors. Independence and the directional second moment bounds imply
\[
\mathbb E_{P^{\otimes n}}
\!\left[\langle\Delta_{1,P},\Delta_{2,P}\rangle^2\right]
\le\Lambda\,\mathbb E_{P^{\otimes n}}\|\Delta_{2,P}\|_2^2
=\Lambda\sum_{j=1}^{J}
\mathbb E_{P^{\otimes n}}
\!\left[\langle\mathbf e^{(j)},\Delta_{2,P}\rangle^2\right]
\le J\Lambda^2.
\]
Consequently,
\[
\operatorname{Var}_{P^{\otimes n}}(W_{\mathrm e})
\le2\Lambda\|\mu_P\|_2^2+J\Lambda^2.
\]
All the required moments are finite because the clipped score and the finite-rank features are bounded.
% lean: CausalSmith.Stat.FinitepHomogeneityDensegamma.oracleBlockVec_quadratic_variance_le

The squared-deviation inequality now gives
\[
\begin{aligned}
&P^{\otimes n}\!\left(
|W_{\mathrm e}-\|\mu_P\|_2^2|
>16\{\sqrt\Lambda\|\mu_P\|_2+\sqrt J\,\Lambda\}
\right)\\
&\quad\le
\frac{2\Lambda\|\mu_P\|_2^2+J\Lambda^2}
{256\{\sqrt\Lambda\|\mu_P\|_2+\sqrt J\,\Lambda\}^2}
\le\frac{2}{256}<0.01.
\end{aligned}
\]
The denominator is positive since \(J\ge1\) and \(\Lambda>0\).
Define the good event by
\[
\mathcal G_P=
\left\{
|W_{\mathrm e}-\|\mu_P\|_2^2|
\le16\{\sqrt\Lambda\|\mu_P\|_2+\sqrt J\,\Lambda\}
\right\}.
\]
Thus \(P^{\otimes n}(\mathcal G_P^c)\le0.01\).
% lean: CausalSmith.Stat.FinitepHomogeneityDensegamma.oracleBlockVec_tail_le

For a null law, on \(\mathcal G_P\) we have
\[
W_{\mathrm e}
\le b^2+16\sqrt\Lambda\,b+16\sqrt J\,\Lambda
\le2b^2+80\sqrt J\,\Lambda
\le2b^2+1024\sqrt J\,\Lambda.
\]
The middle inequality follows from
\[
16\sqrt\Lambda\,b\le b^2+64\Lambda,
\qquad
\Lambda\le\sqrt J\,\Lambda,
\]
the first being equivalent to \((b-8\sqrt\Lambda)^2\ge0\).
The rejection event is therefore contained in \(\mathcal G_P^c\). Since the rule ignores the public randomization variable,
\[
\mathsf R_n^{\mathrm e}(P,\phi^{*,\mathrm e}_{n,v})
=
P^{\otimes n}\!\left(
W_{\mathrm e}>2b^2+1024\sqrt J\,\Lambda
\right)
\le0.01
\qquad(P\in H_0(v)).
\]
These calculations also apply when \(s=1\); when \(J=1\), \(C_J=0\) and both centered block vectors vanish.
% lean: CausalSmith.Stat.FinitepHomogeneityDensegamma.oracle_whole_null_size_le

\item We establish the alternative population bound before applying the same concentration event. Define
\[
\overline\tau_P=\int_0^1\tau_P(x)\,d\lambda(x),
\qquad
c_{P,J}=\frac1{\sqrt J}\sum_{j=1}^{J}(\theta_{\tau,P})_j,
\qquad
h_{P,J}(x)=\tau_P(x)-c_{P,J}.
\]
The coefficient formula for constants and the definition of \(C_J\) give
\[
C_J\theta_{\tau,P}
=\int_0^1h_{P,J}(x)F_J(x)\,d\lambda(x).
\]
On each cell, the histogram projection is the cell average. Hence
\[
\tau_P(x)-(\Pi_J\tau_P)(x)
=J\int_{I_{j,J}}\{\tau_P(x)-\tau_P(z)\}\,d\lambda(z)
\qquad(x\in I_{j,J}),
\]
and \cref{obj:ass:effect-smoothness} gives
\[
|\tau_P(x)-(\Pi_J\tau_P)(x)|\le20J^{-\gamma},
\qquad
\int_0^1\{\tau_P-\Pi_J\tau_P\}^2\,d\lambda
\le(20J^{-\gamma})^2.
\]
Subtracting a constant commutes with \(\Pi_J\). Its cellwise residual has integral zero on every cell, and is therefore orthogonal to the cellwise constant projection. The feature isometry then yields
\[
\int_0^1h_{P,J}^2\,d\lambda
=
\|C_J\theta_{\tau,P}\|_2^2
+\int_0^1\{\tau_P-\Pi_J\tau_P\}^2\,d\lambda.
\]
Expanding the distance to the constant \(c_{P,J}\) also gives
\[
\int_0^1h_{P,J}^2\,d\lambda
=d(P)^2+(\overline\tau_P-c_{P,J})^2
\ge d(P)^2.
\]
It follows that
\[
d(P)^2
\le\|C_J\theta_{\tau,P}\|_2^2+(20J^{-\gamma})^2
\le\bigl(\|C_J\theta_{\tau,P}\|_2+20J^{-\gamma}\bigr)^2,
\]
so
\[
d(P)-20J^{-\gamma}\le\|C_J\theta_{\tau,P}\|_2.
\]
% lean: CausalSmith.Stat.FinitepHomogeneityDensegamma.centered_effect_histogram_distance

The same cellwise coefficient bound used for the null, now applied to \(g_{T,P}-\tau_P\), gives
\[
\|\mu_P-C_J\theta_{\tau,P}\|_2
=
\left\|C_J\int_0^1(g_{T,P}-\tau_P)F_J\,d\lambda\right\|_2
\le b.
\]
Combining this with the preceding lower bound and the reverse triangle inequality yields
\[
d(P)-20J^{-\gamma}-b\le\|\mu_P\|_2.
\]
% lean: CausalSmith.Stat.FinitepHomogeneityDensegamma.oracle_alternative_population_norm_ge

Define the noise scale by
\[
\zeta_{\mathrm{noise}}=\sqrt{\sqrt J\,\Lambda}>0.
\]
If
\[
d(P)\ge20J^{-\gamma}+5b+128\zeta_{\mathrm{noise}},
\]
then
\[
\|\mu_P\|_2\ge4b+128\zeta_{\mathrm{noise}}.
\]
Since \(\sqrt\Lambda\le\zeta_{\mathrm{noise}}\), this implies
\[
16\sqrt\Lambda\,\|\mu_P\|_2\le\frac{\|\mu_P\|_2^2}{2}.
\]
On \(\mathcal G_P\), therefore,
\[
\begin{aligned}
W_{\mathrm e}
&\ge\|\mu_P\|_2^2
-16\sqrt\Lambda\,\|\mu_P\|_2
-16\zeta_{\mathrm{noise}}^2\\
&\ge\frac{\|\mu_P\|_2^2}{2}-16\zeta_{\mathrm{noise}}^2\\
&\ge8b^2+8176\zeta_{\mathrm{noise}}^2
>2b^2+1024\zeta_{\mathrm{noise}}^2.
\end{aligned}
\]
Here the third inequality uses
\[
\frac{(4b+128\zeta_{\mathrm{noise}})^2}{2}
\ge8b^2+8192\zeta_{\mathrm{noise}}^2,
\]
and the last is strict because \(\zeta_{\mathrm{noise}}>0\).
Thus failure to reject is contained in \(\mathcal G_P^c\), proving
\[
1-\mathsf R_n^{\mathrm e}(P,\phi^{*,\mathrm e}_{n,v})\le0.01
\]
under the displayed separation condition.
% lean: CausalSmith.Stat.FinitepHomogeneityDensegamma.oracle_power_of_mean_bound

\item We now bound that sufficient separation by the scale in the statement. Define the variance-growth exponent by
\[
t=\frac{2-p}{p-1}\ge0.
\]
Both dyadic targets are at least one, and upward rounding gives
\[
a^{-1/(p-1)}\le T\le2a^{-1/(p-1)},
\qquad
a^{-1/\gamma}\le J\le2a^{-1/\gamma}.
\]
Since \(1-p<0\) and \(-\gamma<0\),
\[
b=20T^{1-p}\le20a,
\qquad
J^{-\gamma}\le a.
\]
Also \(0\le2-p\le1\), so
\[
T^{2-p}
\le2^{2-p}a^{-t}\le2a^{-t},
\qquad
\sqrt J\le\sqrt2\,a^{-1/(2\gamma)}.
\]
Substitution of the definitions of \(q\) and \(E_0\) gives the balance identity
\[
E_0\left(2+t+\frac1{2\gamma}\right)=1.
\]
Because \(a=s^{-E_0}\), this identity implies
\[
\frac{a^{-t-1/(2\gamma)}}s=a^2.
\]
Consequently,
\[
\sqrt J\,\Lambda
\le160\sqrt2\,\frac{a^{-t-1/(2\gamma)}}s
=160\sqrt2\,a^2,
\qquad
\zeta_{\mathrm{noise}}\le\sqrt{160\sqrt2}\,a.
\]

For every \(n\ge2\), \(n\le3s\). Moreover,
\[
2\gamma q\le2\gamma+q,\qquad 0<E_0\le1,
\qquad 3^{E_0}\le3.
\]
Monotonicity of the negative power therefore gives
\[
a=s^{-E_0}
\le3^{E_0}n^{-E_0}
\le3\rho_n^{\mathrm e}(v).
\]
Combining the approximation, clipping, and covariance bounds yields
\[
\begin{aligned}
20J^{-\gamma}+5b+128\zeta_{\mathrm{noise}}
&\le\left(120+128\sqrt{160\sqrt2}\right)a\\
&\le C^{\mathrm e}_v\rho_n^{\mathrm e}(v).
\end{aligned}
\]
Thus, whenever \(C^{\mathrm e}_v\rho_n^{\mathrm e}(v)<D_v\), every model law at distance at least \(C^{\mathrm e}_v\rho_n^{\mathrm e}(v)\) has type-II error at most \(0.01\).
% lean: CausalSmith.Stat.FinitepHomogeneityDensegamma.oracle_attaining_separation_le

\item Define the upper separation by
\[
\varrho_+=C^{\mathrm e}_v\rho_n^{\mathrm e}(v)>0.
\]
The witnesses already give \(D_v>0\). If \(\varrho_+<D_v\), the calibrated rule satisfies the size constraint \(1/10\), and its alternative errors give
\[
\sup_{P\in\mathcal M_v:\,d(P)\ge\varrho_+}
\left\{1-\mathsf R_n^{\mathrm e}(P,\phi^{*,\mathrm e}_{n,v})\right\}
\le0.01\le\frac1{10}.
\]
For an empty alternative index set, the real-valued supremum is assigned the value \(0\), so the same bound applies in that case. All worst-error suprema are nonnegative, and the admissible test family contains the calibrated rule. Taking its value in the infimum of \cref{obj:def:oracle-risk} therefore shows
\[
B_n^{\mathrm e}(v,\varrho_+)\le\frac1{10}.
\]
Hence \(\varrho_+\) is a candidate in \cref{obj:def:oracle-radius}, and
\[
r_n^{*,\mathrm e}(v)\le\varrho_+.
\]
If instead \(D_v\le\varrho_+\), the candidate \(D_v\) gives
\[
r_n^{*,\mathrm e}(v)\le D_v\le\varrho_+.
\]
Both cases establish the upper radius bound.
% lean: CausalSmith.Stat.FinitepHomogeneityDensegamma.oracleCriticalRadius_le_of_test

\item Retain the prior and null law already constructed, and let \(\phi^{\mathrm e}\) satisfy the stipulated size bound on \(H_0(v)\). Define its seed-averaged map with the common propensity supplied by
\[
\overline\phi^{\mathrm e}_{P_0}(\mathcal D_n)
=\int_0^1\phi^{\mathrm e}(e_{P_0},\mathcal D_n,u_{\mathrm{seed}})\,d\lambda(u_{\mathrm{seed}}),
\qquad
0\le\overline\phi^{\mathrm e}_{P_0}\le1.
\]
This is a measurable function on the complete-record sample space. Integrating event-probability differences over its level sets gives
\[
\left|
\int\overline\phi^{\mathrm e}_{P_0}\,d\mathbb P_{\pi_1,n}
-\int\overline\phi^{\mathrm e}_{P_0}\,dP_0^{\otimes n}
\right|
\le\operatorname{TV}\!\left(P_0^{\otimes n},\mathbb P_{\pi_1,n}\right).
\]
Indeed, each integral equals the integral over \(\vartheta_{\mathrm{lev}}\in[0,1]\) of the probability of the measurable level set
\(\{\overline\phi^{\mathrm e}_{P_0}>\vartheta_{\mathrm{lev}}\}\).

Every positive-weight support law has propensity \(e_{P_0}\). Finite linearity of integration and the null size bound thus yield
\[
\sum_{j=1}^{k}\omega_j
\mathsf R_n^{\mathrm e}(P_j,\phi^{\mathrm e})
\le
\mathsf R_n^{\mathrm e}(P_0,\phi^{\mathrm e})
+\operatorname{TV}\!\left(P_0^{\otimes n},\mathbb P_{\pi_1,n}\right)
<\frac1{10}+\frac12=\frac35.
\]
If every positive-weight index had type-II error below \(2/5\), its rejection probability would exceed \(3/5\). The zero-weight terms vanish, and normalization would then imply
\[
\sum_{j=1}^{k}\omega_j
\mathsf R_n^{\mathrm e}(P_j,\phi^{\mathrm e})
\ge\frac35\sum_{j=1}^{k}\omega_j=\frac35,
\]
a contradiction. Therefore
\[
\exists j:\quad
\omega_j>0,\qquad
1-\mathsf R_n^{\mathrm e}(P_j,\phi^{\mathrm e})\ge\frac25.
\]
This proves the asserted individual error witness for the same common-propensity prior.
% lean: CausalSmith.Stat.FinitepHomogeneityDensegamma.oracle_single_null_prior_error_witness

\item For \(n\ge N^{\mathrm e}_v\), the definition of the ceiling gives
\[
n\ge
\max\left\{2,\left(\frac{2C^{\mathrm e}_v}{d_0}\right)^{1/E_0}\right\}.
\]
Since \(E_0>0\), raising the lower bound for \(n\) to the power \(-E_0\) reverses the inequality:
\[
n^{-E_0}
\le
\left\{
\left(\frac{2C^{\mathrm e}_v}{d_0}\right)^{1/E_0}
\right\}^{-E_0}
=\left(\frac{2C^{\mathrm e}_v}{d_0}\right)^{-1}.
\]
Multiplication by \(C^{\mathrm e}_v>0\) proves
\[
C^{\mathrm e}_v\rho_n^{\mathrm e}(v)\le\frac{d_0}{2}.
\]
% lean: CausalSmith.Stat.FinitepHomogeneityDensegamma.oracle_attainment_threshold

\item Let \(P\in\mathcal M_v\). On the common almost-everywhere set where the two conditional arm-mean identities hold,
\[
\int_{\mathbb R}y\,Q_{0,P}(dy\mid x)=m_{0,P}(x),
\qquad
\int_{\mathbb R}y\,Q_{1,P}(dy\mid x)=m_{0,P}(x)+\tau_P(x).
\]
Overlap gives \(e_P(x)>0\) and \(1-e_P(x)>0\) at every \(x\). Pulling the constant denominators outside the arm integrals and cancelling them therefore yields
\[
\begin{aligned}
&e_P(x)\int_{\mathbb R}\frac{y}{e_P(x)}\,Q_{1,P}(dy\mid x)
+(1-e_P(x))
\int_{\mathbb R}\frac{-y}{1-e_P(x)}\,Q_{0,P}(dy\mid x)\\
&\qquad=
\int_{\mathbb R}y\,Q_{1,P}(dy\mid x)
-\int_{\mathbb R}y\,Q_{0,P}(dy\mid x)
=\tau_P(x).
\end{aligned}
\]
This is the claimed identity for the untruncated score.
% lean: CausalSmith.Stat.FinitepHomogeneityDensegamma.oracle_untruncated_mean_identity

\item Finally, let \(\mathcal K\subseteq\mathcal V\) be compact. The coordinate functions \(p\) and \(\gamma\) are continuous in the inherited topology. Since \(p>1\) and \(2\gamma+(p-1)/p>0\), the function
\[
v\longmapsto
E_0(v)=
\frac{2\gamma\,((p-1)/p)}
{2\gamma+(p-1)/p}
\]
is continuous and positive throughout \(\mathcal V\). If \(\mathcal K\) is nonempty, it attains a positive minimum; if \(\mathcal K\) is empty, choose the value one. Thus there is
\[
\eta_{\mathcal K}>0,
\qquad
\eta_{\mathcal K}\le E_0(v)\quad(v\in\mathcal K).
\]
The multiplier bounds already proved give
\[
c^{\mathrm e}_v\ge\frac1{64\sqrt3},
\qquad
C^{\mathrm e}_v=3\left(120+128\sqrt{160\sqrt2}\right)
\qquad(v\in\mathcal K).
\]
Using \(d_0=1/(8\sqrt2)\) from \cref{obj:lem:paired-tent-full-record-lower}, the base \(2C^{\mathrm e}_v/d_0\) exceeds one and is independent of \(v\). Therefore
\[
N^{\mathrm e}_v
\le
\left\lceil
\max\left\{
2,\left(\frac{2C^{\mathrm e}_v}{d_0}\right)^{1/\eta_{\mathcal K}}
\right\}
\right\rceil
\qquad(v\in\mathcal K).
\]
This supplies a common finite threshold bound. Positivity of the denominators and continuity of the exponent also hold at \(p=2\) and \(\gamma=1/4\), so the same compact-uniform argument includes those faces.
% lean: CausalSmith.Stat.FinitepHomogeneityDensegamma.oracle_compact_uniformity
\end{enumerate}
\end{proof}
